\documentclass[11pt]{article}
\usepackage[T1]{fontenc}
\usepackage{amsmath,amssymb,amsthm}
\usepackage[margin=1in]{geometry}
\usepackage[hidelinks]{hyperref}
\hypersetup{
  pdftitle={Exact Toughness under Tetrahedral Replacement in Plane Triangulations},
  pdfauthor={Xiaojian Zeng}
}

\newtheorem{theorem}{Theorem}
\newtheorem{lemma}[theorem]{Lemma}

\newcommand{\comp}{c}
\newcommand{\tough}{\tau}

\title{Exact Toughness under Tetrahedral Replacement in Plane Triangulations}
\author{Xiaojian Zeng\\
PBC School of Finance, Tsinghua University, 100190\\
\texttt{xiaojianzeng@foxmail.com}}
\date{}

\begin{document}
\maketitle

\begin{abstract}
Let \(G\) be a finite plane triangulation and let \(v\) be a vertex of
degree three.  The tetrahedral replacement \(R_v(G)\) deletes \(v\), inserts
a copy of \(K_4\), and attaches three of its vertices cyclically to the three
neighbors of \(v\).  We prove the exact identity
\[
  \tough(R_v(G))=\min\{\tough(G),3/2\},
\]
with the convention \(\tough(K_n)=\infty\).  The proof transfers every
toughness level at most \(3/2\) through the replacement and, in the reverse
direction, shows that an optimum toughness cut can avoid the replaced
vertex.  The complete boundary case \(G=K_4\) is handled separately.
\end{abstract}

\section{Statement and context}

For a graph \(J\), write \(\comp(J)\) for its number of connected components.
Following Chv\'{a}tal~\cite{Chvatal1973}, a graph is \(t\)-tough if
\[
  |S|\ge t\,\comp(J-S)
  \quad\text{whenever }S\subseteq V(J)\text{ and }\comp(J-S)\ge2.
\]
For a finite noncomplete graph this is equivalent to
\[
  \tough(J)=
  \min_{\substack{S\subseteq V(J)\\\comp(J-S)\ge2}}
       \frac{|S|}{\comp(J-S)},
\]
and we use the standard convention \(\tough(K_n)=\infty\).

Let \(v\) have neighbors \(a,b,c\), listed in their cyclic order.  Since
\(G\) is a plane triangulation, these neighbors induce a triangle.  Put
\[
 B=\{a,b,c\},\qquad A=\{x,y,z,w\}.
\]
The graph \(H=R_v(G)\) has vertex set
\((V(G)\setminus\{v\})\mathbin{\dot\cup}A\).  Its edges are those of
\(G-v\), all six edges within \(A\), and
\[
 xa,\ xb,\ yb,\ yc,\ zc,\ za;
\]
there are no other edges incident with \(A\).  Shan~\cite{Shan2026}
introduced this operation and proved in Theorem~1.4(ii) that
\(3/2\)-toughness is preserved.  The result below gives the exact value,
including inputs of toughness below \(3/2\).  Thus the additional content
relative to that cited version is a parameterized lower transfer together
with a matching reverse transfer; no broader priority claim is intended.

\begin{theorem}\label{thm:main}
For every finite plane triangulation \(G\) and every degree-three vertex
\(v\in V(G)\),
\[
  \tough(R_v(G))=\min\{\tough(G),3/2\}.
\]
\end{theorem}

\section{The two transfers}

\begin{lemma}[Lower transfer]\label{lem:lower}
If \(0\le t\le3/2\) and \(G\) is \(t\)-tough, then \(R_v(G)\) is
\(t\)-tough.
\end{lemma}

\begin{proof}
Let \(H=R_v(G)\), and fix \(S\subseteq V(H)\) for which
\(d:=\comp(H-S)\ge2\).  Set
\[
 S_0=S\setminus A,\qquad k=|S\cap A|,
 \qquad g=\comp(G-S_0),\qquad
 g'=\comp((G-v)-S_0).
\]
All surviving vertices of \(B\), if there are any, lie in one component of
\((G-v)-S_0\), because \(B\) is a clique.  Adding \(v\) therefore gives
\begin{equation}\label{eq:g-gprime}
 g=\begin{cases}
      g',&B\setminus S_0\ne\varnothing,\\
      g'+1,&B\subseteq S_0.
    \end{cases}
\end{equation}
We consider all possible values of \(k\).

If \(k=0\) or \(k=1\), then \(H[A\setminus S]\) is nonempty and connected.
Moreover, every vertex of \(B\setminus S_0\) retains a neighbor in
\(A\setminus S\), since each vertex of \(B\) has two neighbors among
\(x,y,z\).  Relative to the common graph \((G-v)-S_0\), adding
\(A\setminus S\) has exactly the same effect on the component count as
adding \(v\).  Hence \(d=g\ge2\), and the \(t\)-toughness of \(G\) yields
\[
  |S|\ge |S_0|\ge tg=td.
\]

Suppose next that \(k=2\) or \(k=3\).  The surviving part of \(A\) is again
nonempty and connected.  Adding it to \((G-v)-S_0\) can increase the
component count by at most one, so by~\eqref{eq:g-gprime},
\begin{equation}\label{eq:component-upper}
  d\le g'+1\le g+1.
\end{equation}
If \(g\ge2\), then
\[
 |S|=|S_0|+k\ge tg+k\ge t(g+1)\ge td,
\]
where \(k\ge2\ge t\).  If \(g=1\), admissibility and
\eqref{eq:component-upper} force \(d=2\).  We then have
\(|S|\ge3\ge2t=td\), unless \(k=2\) and \(S_0=\varnothing\).  That exception
cannot disconnect \(H\).  Indeed, \(G-v\) is connected: every component of
\(G-v\) contains a neighbor of \(v\), and the three such neighbors form a
clique.  The two surviving gadget vertices are adjacent, and at least one is
among \(x,y,z\), hence is adjacent to \(B\).  Thus \(H-S\) would be connected,
contrary to \(d=2\).

Finally, if \(k=4\), then
\[
  H-S=(G-v)-S_0=G-(S_0\cup\{v\}).
\]
This graph has \(d\ge2\) components, so \(t\)-toughness applied to the cut
\(S_0\cup\{v\}\) gives \(|S_0|+1\ge td\).  Consequently
\(|S|=|S_0|+4\ge td\).  All admissible cuts of \(H\) have now been covered.
\end{proof}

\begin{lemma}[An optimum avoids \(v\)]\label{lem:avoid}
Let \(G\) be finite and noncomplete.  If \(v\) has a clique as its
neighborhood, then a toughness-minimizing set of \(G\) avoids \(v\).
\end{lemma}

\begin{proof}
Because \(G\) is finite and noncomplete, an admissible minimizing set \(Q\)
exists.  Write \(d=\comp(G-Q)\ge2\), and suppose that \(v\in Q\).  Put
\(Q_0=Q\setminus\{v\}\).

If \(N_G(v)\nsubseteq Q\), all surviving neighbors of \(v\) belong to one
component, since they form a clique.  Restoring \(v\) therefore merges no
components, and \(\comp(G-Q_0)=d\).  The numerator has decreased by one,
contradicting minimality.

If \(N_G(v)\subseteq Q\), restoring \(v\) creates one isolated component,
so \(\comp(G-Q_0)=d+1\).  The set \(Q_0\) is admissible and
\[
  \frac{|Q|-1}{d+1}<\frac{|Q|}{d},
\]
again a contradiction.  Hence a minimizing set avoids \(v\).
\end{proof}

\section{Proof of the exact law}

\begin{proof}[Proof of Theorem~\ref{thm:main}]
First assume \(G\ne K_4\).  A plane triangulation with a degree-three vertex
has at least four vertices, and the order-four case is \(K_4\); hence
\(|V(G)|\ge5\).  Deleting \(B=N_G(v)\) leaves \(v\) isolated and leaves at
least one further vertex.  Thus \(G\) is noncomplete and
\begin{equation}\label{eq:upper-three-halves}
  \tough(G)\le \frac{|B|}{\comp(G-B)}\le\frac32.
\end{equation}
Apply Lemma~\ref{lem:lower} with \(t=\tough(G)\) to obtain
\(\tough(H)\ge\tough(G)\).

For the other inequality, Lemma~\ref{lem:avoid} supplies a minimizing set
\(Q\subseteq V(G)\setminus\{v\}\).  Relative to the common graph
\((G-v)-Q\), the vertex \(v\) in \(G-Q\) and the connected gadget \(A\) in
\(H-Q\) have the same effect.  If \(B\setminus Q\ne\varnothing\), each
attaches to the component containing the surviving part of \(B\); if
\(B\subseteq Q\), each forms one isolated component.  Therefore
\[
  \comp(H-Q)=\comp(G-Q),
\]
and
\[
  \tough(H)\le\frac{|Q|}{\comp(H-Q)}
  =\frac{|Q|}{\comp(G-Q)}=\tough(G).
\]
Together with~\eqref{eq:upper-three-halves}, this proves the claimed identity
when \(G\ne K_4\).

Now take \(G=K_4\).  The convention gives \(\tough(G)=\infty\), while
Lemma~\ref{lem:lower} with \(t=3/2\) gives \(\tough(H)\ge3/2\).  In \(H\),
deleting \(\{x,y,z\}\) leaves exactly two components: the isolated vertex
\(w\) and the triangle on \(B\).  Hence \(\tough(H)\le3/2\), proving
\[
  \tough(H)=\frac32=\min\{\infty,3/2\}.
\]
\end{proof}

\section*{Computational and AI assistance}

The theorem is established by the preceding argument; no computation is
used in its proof.  As a supplementary falsification check, a Python program
using exact rational arithmetic enumerated every graph-atlas plane
triangulation through seven vertices having a degree-three vertex, all
sixteen ordered choices of \((G,v)\), and every vertex subset in both \(G\)
and \(R_v(G)\).  All checks agreed with Theorem~\ref{thm:main}, including the
exact \(K_4\) value \(3/2\).  For comparison,
Zeng--Liu--Ratnavelu--Huat~\cite{ZengEtAl2026Certificates} likewise pair
exhaustive exact finite enumeration with independent exact verification while
explicitly leaving the unrestricted ambient problems unresolved.  Generative AI
assisted with organizing the proof exposition, checking the case split, and
screening bibliographic metadata.  The author reviewed the mathematical
argument and is responsible for the manuscript.

\bibliographystyle{plain}
\bibliography{references}

\end{document}
